\textbf{Proof.}
Fix \(\bar p\in(0,1)\) and let \(n\in\mathcal N_{\bar p}\) be arbitrary. Write a no-interference sign schedule as \(\alpha=(b_{i0},b_{i1})_{i\in V_n}\). For schedules \(\alpha\) and \(\beta\), let
\[
 D_0=\{i:b_{i0}^\alpha\ne b_{i0}^\beta\},
 \qquad
 D_1=\{i:b_{i1}^\alpha\ne b_{i1}^\beta\}. \tag{4}
\]
They collide at assignment \(z\) exactly when every member of \(D_0\) is treated and every member of \(D_1\) is controlled. If \(D_0\cap D_1\ne\varnothing\), collision is impossible. Otherwise put \(a=|D_0|\) and \(b=|D_1|\). Under the uniform law on \(\Omega_n\), the collision probability is
\[
 \frac{\binom{n-a-b}{m_n-a}}{\binom n{m_n}}, \tag{5}
\]
interpreted as zero unless \(a\le m_n\) and \(b\le n-m_n\). Also,
\[
 |t_\alpha-t_\beta|
 \le\frac{2(a+b)}n. \tag{6}
\]
For every disjoint pair of sets of sizes \((a,b)\), equality in \((6)\) is attained by orienting all treatment-response disagreements in one direction and all control-response disagreements in the opposite direction. Thus \((5)\)--\((6)\) give the maximum in \((1)\) under complete randomization.

It remains to verify its optimality over designs. Every permutation of the units acts equivariantly on assignments and schedules, preserves targets, and relabels observed sign vectors bijectively. For any design \(\pi\), average its images under all unit permutations. The resulting law is uniform on \(\Omega_n\). Since \(C_n(\pi)^2\) is a maximum of linear functions of \(\pi\), it is convex, while equivariance preserves its value. Jensen's inequality therefore gives
\[
 C_n(\pi^{\mathrm{unif}})^2
 \le \frac1{n!}\sum_g C_n(g\pi)^2
 =C_n(\pi)^2. \tag{7}
\]
This proves \((1)\).

Fix an assignment and a compatible sign observation. At every unit exactly one of the two endpoint signs is observed and the other remains free. Varying that unobserved sign changes the target range by \(2/n\), independently across the \(n\) units. The target-fiber diameter is therefore exactly two, which proves \((2)\).

Put \(p_n=m_n/n\) and \(q_n=\max\{p_n,1-p_n\}\). Taking either \((a,b)=(1,0)\) or \((0,1)\) in \((1)\) gives
\[
 \inf_\pi C_n(\pi)\ge\frac{2\sqrt{q_n}}n. \tag{8}
\]
For \(k=a+b\), the probability in \((5)\) is at most \(p_n^a\) and at most \((1-p_n)^b\). Since \(\max\{a,b\}\ge k/2\), it is at most \(q_n^{k/2}\). Choose a constant \(q\) with
\[
 \max\{\bar p,1-\bar p\}<q<1.
\]
Eventually \(q_n\le q\), and \((1)\) then yields
\[
 \inf_\pi C_n(\pi)
 \le\frac2n\sup_{k\ge1}kq^{k/4}
 =\frac{2M(q)}n,
 \qquad M(q)<\infty. \tag{9}
\]
Equations \((8)\)--\((9)\) prove \((3)\) and refute every size-uniform reverse constant.

For the orbit claim, an ordered schedule pair assigns one of \(16\) four-sign types \((b_{i0}^\alpha,b_{i1}^\alpha,b_{i0}^\beta,b_{i1}^\beta)\) to each unit. Its permutation orbit is determined by the counts of these types, so there are at most \(\binom{n+15}{15}\) pair orbits; passage to unordered pairs cannot increase this number. There is one exact-count assignment orbit. Formula \((5)\) shows that the orbit coefficients are ratios of binomial coefficients and have polynomial bit length.

Finally, take any equivariant fractional diameter cover with congestion \(\kappa_n\). The uniform design is invariant under every action on \(\Omega_n\). Summing the cover inequality against that design and bounding each weighted pairwise collision term by \(C_n(\pi^{\mathrm{unif}})^2\) gives
\[
 4=U_n(\pi^{\mathrm{unif}})^2
 \le\kappa_n C_n(\pi^{\mathrm{unif}})^2
 \le\frac{4M(q)^2}{n^2}\kappa_n. \tag{10}
\]
Thus \(\kappa_n\ge n^2/M(q)^2\) eventually. Radius-zero locality follows directly from \(e^{\mathrm{NI}}_{i,n}(z)=z_i\). No positivity beyond the nonempty exact-budget slice is used. \(\square\)